% !PRINTED 482
\weberpage{482}
%<<CONT
\weq{1}{\alpha, \alpha_1, \alpha_2 \ldots \alpha_{n-1}}
and the remaining roots of $F(x)$ be
\weq{2}{\alpha_n, \alpha_{n+1} \ldots \alpha_{m-1}.}

If $\alpha'$ is any one of the quantities (1), and $\alpha''$ one of the
quantities (2), then no permutation can occur in the group $P$ by which
$\alpha'$ is carried into $\alpha''$; for by hypothesis
\[
f(\alpha') = 0,
\]
and if now a permutation occurred in $P$ by which $\alpha'$ were replaced by
$\alpha''$, then on account of the property \S.\,149, a) we should also have
$f(\alpha'') = 0$, which contradicts our assumption.

\gesperrt{Thus by the permutations of} $P$ \gesperrt{the roots (1) of} $f(x)$
\gesperrt{are permuted only among themselves.}

If conversely the group $P$ has the property that its permutations permute a
part of the $m$ roots $\alpha$, such as the system (1), only among
themselves, then the product
\[
(x - \alpha)(x - \alpha_1) \ldots (x - \alpha_{n-1}) = f(x)
\]
admits all the permutations of $P$, and is therefore by \S.\,149, b)
contained in $\Omega$; that is, $F(x)$ is reducible.

A permutation group is called \gesperrt{transitive} if it contains at least
one permutation which carries an arbitrary element into any other arbitrary
element; in the opposite case, that is, when the elements can be so
decomposed into two or more portions that by no permutation of the group
does an element of the one portion pass into an element of the other,
the group is called
\gesperrt{intransitive}. Accordingly we may sum up what has been proved in
the theorem:

% [sic] The printed theorem reverses the correspondence proved immediately above: irreducible corresponds to transitive, and reducible to intransitive. The source wording is preserved.
\gesperrt{1. The equation $F(x) = 0$ is reducible or irreducible, according
as its Galois group is transitive or intransitive.}

If a part of the elements is so connected that by permutations from $P$
every element of this part can pass into every other, we call the elements
of this part \gesperrt{transitively connected}. The various transitively
connected systems, which by all the permutations of $P$ are permuted only
among themselves, are called the \gesperrt{systems of
intransitivity}.
